\chapter{مفاهیم پایه و کارهای مرتبط}

\section{مقدمه}

در این فصل، مفاهیم و پژوهش‌های مرتبط با تولید و ارزیابی توضیحات شطرنجی مبتنی بر مدل‌های زبانی بررسی می‌شوند. مسئله مورد بررسی در این پروژه در نقطه تلاقی چند حوزه قرار دارد: تحلیل شطرنج توسط موتورهای تخصصی، مدل‌های زبانی و تولید زبان طبیعی، تولید توضیحات برای تصمیمات پیچیده و ارزیابی کیفیت خروجی مدل‌های زبانی. بنابراین، برای درک روش پیشنهادی، لازم است ابتدا ارتباط میان این مفاهیم مشخص شود.

در ساده‌ترین حالت، می‌توان فرایند مورد نظر را به این صورت در نظر گرفت که یک موقعیت شطرنجی به‌عنوان ورودی دریافت می‌شود و یک موتور شطرنج با بررسی وضعیت صفحه، اطلاعاتی درباره کیفیت حرکت و تغییرات موقعیت ارائه می‌کند. این اطلاعات در مرحله بعد می‌تواند به یک مدل زبانی داده شود تا آن را به توضیحی قابل فهم برای انسان تبدیل کند. در نهایت، توضیح تولید شده باید با استفاده از معیارهای مشخص ارزیابی شود. بنابراین، موتور شطرنج نقش تأمین اطلاعات تخصصی و قابل بررسی را دارد، مدل زبانی وظیفه تبدیل این اطلاعات به زبان طبیعی را بر عهده می‌گیرد و بخش ارزیابی کیفیت خروجی را بررسی می‌کند.

یکی از دلایل اهمیت این زنجیره آن است که هر بخش محدودیت‌های خاص خود را دارد. موتورهای شطرنج در ارزیابی موقعیت و انتخاب حرکت بسیار قدرتمند هستند، اما خروجی آن‌ها معمولاً به‌صورت عددی یا در قالب خطوط بازی ارائه می‌شود و برای آموزش مستقیم انسان مناسب نیست. در مقابل، مدل‌های زبانی در تولید متن روان و قابل فهم توانایی بالایی دارند، اما ممکن است در توضیح موقعیت‌های تخصصی، اطلاعات نادرست تولید کنند. پژوهش‌های جدید نیز این فاصله میان دقت مدل‌های تخصصی و قابلیت توضیح‌دهی مدل‌های زبانی را به‌عنوان یکی از چالش‌های اصلی این حوزه مطرح کرده‌اند.

\subsection{تحلیل شطرنج و موتورهای شطرنج}

موتور شطرنج نرم‌افزاری است که با بررسی وضعیت صفحه و جست‌وجوی حرکت‌های ممکن، تالش می‌کند بهترین تصمیم را برای یک موقعیت مشخص پیدا کند. خروجی موتور می‌تواند شامل حرکت پیشنهادی، ارزیابی موقعیت و خطوط احتمالی ادامه بازی باشد. در پروژه حاضر، موتور  \lr{Stockfish}به‌عنوان منبع تحلیل تخصصی مورد استفاده قرار می‌گیرد.

اهمیت موتور شطرنج در این پروژه تنها در پیدا کردن بهترین حرکت نیست. اطلاعات موتور به‌عنوان یک مبنای قابل بررسی برای تشخیص تغییرات ایجادشده در نتیجه یک حرکت نیز استفاده می‌شود. به این ترتیب، مدل زبانی می‌تواند به جای تولید توضیح صرفاً بر اساس الگوهای زبانی، از اطلاعات حاصل از تحلیل موقعیت نیز استفاده کند.

این ایده با برخی پژوهش‌های پیشین نیز هم‌راستا است. برای نمونه \lr{Lee}، و همکاران نشان داده‌اند که ترکیب مدل‌های زبانی با موتورهای استدالل نمادین می‌تواند برای تولید توضیحات شطرنجی مفید باشد و این ترکیب، امکان استفاده هم‌زمان از توانایی تحلیلی موتور و توانایی زبانی مدل را فراهم می‌کند.

\subsection{مدل‌های زبانی و تولید متن}

مدل زبانی، مدلی یادگیری‌محور است که با استفاده از الگوهای موجود در داده‌های متنی می‌تواند متن تولید کند، دستورها را دنبال کند و اطلاعات ورودی را به زبان طبیعی بازنمایی کند. مدل‌های زبانی جدید در بسیاری از وظایف تولید متن عملکرد قابل توجهی دارند و می‌توانند برای تبدیل داده‌های ساختاریافته به توضیحات قابل فهم استفاده شوند.

در مسئله شطرنج، استفاده از مدل زبانی به این دلیل اهمیت دارد که خروجی موتور شطرنج برای بسیاری از کاربران، به‌خصوص بازیکنان مبتدی، مستقیما ً قابل تفسیر نیست. برای مثال، دانستن اینکه ارزیابی یک موقعیت پس از یک حرکت کاهش یافته است، به‌تنهایی توضیح نمی‌دهد که این کاهش به دلیل ضعف امنیت شاه، از دست دادن یک مهره، ایجاد ضعف در ساختار پیاده‌ها یا از دست دادن یک فرصت تاکتیکی رخ داده است.

با این حال، مدل زبانی به‌تنهایی تضمین نمی‌کند که توضیح تولید شده از نظر شطرنجی صحیح باشد. پژوهش‌های مرتبط نشان داده‌اند که مدل‌های زبانی می‌توانند توضیحاتی روان تولید کنند، در حالی که برخی ادعاهای موجود در آن‌ها با وضعیت واقعی صفحه سازگار نیست. به همین دلیل، استفاده از اطلاعات موتور شطرنج و روش‌های ارزیابی مستقل اهمیت پیدا می‌کند.

\subsection{تولید توضیحات شطرنجی}

تولید توضیحات شطرنجی را می‌توان به‌عنوان مسئله‌ای برای تبدیل اطلاعات یک بازی یا موقعیت شطرنجی به متن طبیعی تعریف کرد. در پژوهش  \lr{Jhamtani}~\cite{jhamtani2018}و همکاران، این مسئله با استفاده از مجموعه‌داده بزرگی شامل بیش از  ۲۹۸هزار جفت حرکت و توضیح در بیش از  ۱۱هزار بازی بررسی شده است. این پژوهش نشان می‌دهد که توضیح دادن حرکات شطرنجی تنها به وضعیت فعلی صفحه وابسته نیست و ویژگی‌های مختلف بازی و زمینه حرکت نیز می‌توانند در نوع توضیح تولیدشده مؤثر باشند.

در پژوهش  \lr{Zang}~\cite{zang2019}و همکاران نیز تولید خودکار توضیحات شطرنجی در دسته‌هایی مانند توصیف، مقایسه و برنامه‌ریزی بررسی شده است. این پژوهش با وارد کردن یک موتور شطرنج عصبی در مدل تولید متن، تالش می‌کند اطلاعات مربوط به صفحه و تصمیم شطرنجی را با فرایند تولید زبان طبیعی ترکیب کند. ارزیابی این روش نیز هم با معیارهای خودکار و هم با ارزیابی انسانی انجام شده است.

این مطالعات نشان می‌دهند که تولید توضیح شطرنجی یک مسئله صرفاً زبانی نیست و برای تولید توضیح مناسب، درک وضعیت صفحه و تصمیم اتخاذشده نیز اهمیت دارد.

\subsection{اتصال تحلیل تخصصی و مدل زبانی}

یکی از رویکردهای مهم در پژوهش‌های جدید، ترکیب یک سامانه تخصصی با مدل زبانی است. در این رویکرد، سامانه تخصصی اطلاعات دقیق مربوط به مسئله را فراهم می‌کند و مدل زبانی وظیفه تبدیل آن اطلاعات به یک خروجی قابل فهم را بر عهده می‌گیرد.

دکتر  \lr{Kim}~\cite{kim2025}و همکاران در سال  ۲۰۲۵با معرفی روش \lr{Concept-guided Chess Commentary Generation}تالش کردند توانایی تصمیم‌گیری مدل‌های تخصصی را با قابلیت تولید زبان طبیعی مدل‌های زبانی ترکیب کنند. در این روش، توضیحات بر اساس مفاهیم شطرنجی هدایت می‌شوند و در کنار روش تولید، یک روش ارزیابی مبتنی بر دانش تخصصی نیز معرفی شده است. نتایج این پژوهش با ارزیابی انسانی و ارزیابی مبتنی بر مدل زبانی بررسی شده‌اند.

این رویکرد به مسئله پروژه حاضر نزدیک است؛ زیرا در هر دو، هدف آن است که اطلاعات تخصصی شطرنج به توضیح قابل فهم تبدیل شود. تفاوت مهم پروژه حاضر در تمرکز آن بر ایجاد یک مجموعه آزمون متوازن، مقایسه تعداد زیادی مدل زبانی مختلف و ارزیابی خروجی آن‌ها در یک فرایند یکسان است.

\subsection{ارزیابی توضیحات شطرنجی}

یکی از چالش‌های مهم این حوزه، تعیین کیفیت توضیح تولید شده است. ارزیابی یک متن شطرنجی تنها بر اساس روان بودن آن کافی نیست؛ زیرا ممکن است متن از نظر زبانی مناسب باشد اما ادعاهای شطرنجی نادرستی داشته باشد.

پژوهش‌های اولیه از ارزیابی انسانی و معیارهای خودکار برای بررسی کیفیت توضیحات استفاده کرده‌اند. برای مثال، \lr{Jhamtani}~\cite{jhamtani2018}و همکاران عالوه بر مقایسه مدل‌ها، بخشی از خروجی‌ها را توسط انسان‌ها از نظر صحت و روانی بررسی کردند.

در پژوهش  \lr{Kim}~\cite{kim2025}و همکاران نیز عالوه بر تولید توضیح، یک چارچوب ارزیابی مبتنی بر مدل زبانی و دانش تخصصی معرفی شده است. این موضوع نشان می‌دهد که مسئله ارزیابی توضیحات به اندازه مسئله تولید آن اهمیت دارد.

پژوهش جدیدتر  \lr{Hebbar}~\cite{hebbar2026}و همکاران در سال  ۲۰۲۶این مسئله را با تمرکز مستقیم بر توهم در توضیحات شطرنجی بررسی می‌کند. در این پژوهش، چارچوب  \lr{ACT-Eval}ادعاهای موجود در توضیحات را به اجزای کوچک‌تر تقسیم کرده و آن‌ها را با استفاده از ابزارهای مبتنی بر موتور و مراجع تأییدشده بررسی می‌کند. این پژوهش همچنین یک بنچمارک شامل  ۳۲۵جفت موقعیت-حرکت ارائه می‌دهد و نشان می‌دهد که خطاهای محتوایی در توضیحات مدل‌های زبانی همچنان مسئله‌ای جدی هستند.

پژوهش  \lr{Cui}~\cite{cui2026}و همکاران نیز به‌طور مستقیم مسئله بیان استراتژی‌های شطرنج در زبان طبیعی را بررسی کرده و عالوه بر روش تولید توضیح، یک چارچوب برای ارزیابی عینی توضیحات پیشنهاد می‌کند. این پژوهش همچنین محدودیت ارزیابی صرفاً مبتنی بر مدل‌های زبانی را مورد توجه قرار می‌دهد.

در نتیجه، در پروژه حاضر نیز ارزیابی خروجی‌ها به‌عنوان یکی از اجزای اصلی در نظر گرفته می‌شود. برای این منظور، خروجی مدل‌های مختلف با استفاده از معیارهای یکسان و توسط یک مدل ارزیاب مستقل بررسی می‌شوند.

\section{تحلیل نقاط قوت و ضعف منابع غیرپژوهشی مشابه}

در کنار پژوهش‌های دانشگاهی، برخی سامانه‌های تجاری و عمومی شطرنج نیز قابلیت‌هایی برای توضیح متنی حرکات و موقعیت‌های شطرنجی ارائه می‌کنند. هدف اصلی این سامانه‌ها معمولاً تحلیل و آموزش شطرنج برای کاربران است و نه بررسی و مقایسه مدل‌های زبانی. بااین‌حال، از آنجا که برخی از این سامانه‌ها برای حرکات شطرنجی توضیحات متنی تولید می‌کنند، می‌توان آن‌ها را از نظر قابلیت‌های ارائه‌شده و تفاوت با سامانه پیشنهادی این پروژه بررسی کرد.

یکی از نقاط قوت این سامانه‌ها، ارائه یک محیط یکپارچه و آماده استفاده برای کاربر است. کاربر می‌تواند موقعیت یا بازی شطرنجی خود را در اختیار سامانه قرار دهد، تحلیل شطرنج را مشاهده کند و در برخی موارد، عالوه بر اطلاعات مربوط به حرکت، یک توضیح متنی نیز دریافت کند. این موضوع باعث می‌شود استفاده از قابلیت‌های تحلیلی برای کاربران، به‌ویژه کاربران غیرمتخصص، ساده‌تر شود.

بااین‌حال، نحوه تولید توضیحات متنی در بسیاری از این سامانه‌ها برای کاربر قابل مشاهده نیست. به‌طور مشخص، مشخص نیست که متن تولیدشده دقیقاً توسط چه مدل یا ترکیبی از چه روش‌هایی ایجاد می‌شود و چه اطلاعاتی از وضعیت شطرنج برای تولید آن در اختیار سامانه قرار می‌گیرد. در نتیجه، امکان بررسی دقیق ارتباط میان تحلیل شطرنج، اطلاعات ورودی و متن نهایی تولیدشده برای یک کاربر یا پژوهشگر وجود ندارد.

محدودیت دیگر، وابستگی این قابلیت‌ها به سرویس‌های آنالین است. توضیحات متنی در این سامانه‌ها معمولاً در سمت سرور تولید می‌شوند و کاربر تنها نتیجه نهایی را دریافت می‌کند. بنابراین، مدل یا فرایند تولیدکننده متن لزوماً به‌صورت محلی در اختیار کاربر قرار نمی‌گیرد. این مسئله با هدف پروژه حاضر که بررسی امکان استفاده از مدل‌های زبانی به‌صورت محلی و قابل کنترل است، تفاوت دارد.

از طرف دیگر، این سامانه‌ها برای مقایسه چند مدل زبانی مختلف نیز طراحی نشده‌اند. در پروژه حاضر، هدف آن است که مدل‌های مختلف با ورودی یکسان روی یک مجموعه آزمون مشخص آزمایش شوند و خروجی آن‌ها با استفاده از معیارهای یکسان مورد ارزیابی قرار گیرد. به همین دلیل، تمرکز پروژه حاضر نه بر رقابت با امکانات عمومی سامانه‌های شطرنجی، بلکه بر فراهم کردن یک محیط قابل بررسی برای مطالعه و مقایسه تولید توضیحات شطرنجی توسط مدل‌های زبانی است.

تفاوت اصلی میان سامانه‌های مشابه و پروژه حاضر در جدول~\ref{tab:related-systems} خلاصه شده است.

\begin{table}[htbp]
\centering
\caption{مقایسه قابلیت‌های سامانه‌های عمومی شطرنج با پروژه حاضر}
\label{tab:related-systems}
\begingroup
\small\setstretch{1.2}\renewcommand{\arraystretch}{1.45}
\setlength{\tabcolsep}{7pt}
\rowcolors{2}{gray!7}{white}
\begin{tabularx}{\linewidth}{>{\raggedleft\arraybackslash}p{3.5cm} >{\raggedleft\arraybackslash}X >{\raggedleft\arraybackslash}X}
\toprule
\rowcolor{gray!22}
\textbf{ویژگی} & \textbf{سامانه‌های تجاری و عمومی} & \textbf{پروژه حاضر} \\
\midrule
هدف اصلی & تحلیل و آموزش شطرنج & تولید و ارزیابی توضیحات شطرنجی \\
تحلیل موقعیت شطرنجی & معمولاً وجود دارد & وجود دارد \\
تولید توضیح متنی & در برخی سامانه‌ها & دارد \\
شفافیت فرایند تولید & معمولاً محدود یا نامشخص & قابل بررسی و کنترل \\
دسترسی به مدل مولد & معمولاً در اختیار کاربر نیست & امکان استفاده از مدل محلی \\
اجرای کامل محلی & معمولاً فراهم نیست & با انتخاب مدل محلی امکان‌پذیر است \\
آزمایش چند مدل & هدف اصلی سامانه نیست & امکان مقایسه مدل‌های مختلف \\
ارزیابی مستقل خروجی & معمولاً ارائه نمی‌شود & داوری مشترک برای پاسخ مدل‌ها \\
\bottomrule
\end{tabularx}
\endgroup
\par\smallskip
{\footnotesize این جدول، جمع‌بندی کیفی بحث این بخش است؛ امکانات سرویس‌های عمومی یکسان نیست و جدول رتبه‌بندی تجربی آن‌ها محسوب نمی‌شود.}
\end{table}


بنابراین، سامانه‌های تجاری را نمی‌توان مستقیما ً جایگزین یک چارچوب ارزیابی پژوهشی دانست. ارزش آن‌ها در ارائه یک تجربه کاربری آماده و قابلیت‌های عملی تحلیل شطرنج است، در حالی که پروژه حاضر بر قابل بررسی بودن فرایند تولید، اجرای محلی، آزمایش مدل‌های مختلف و ارزیابی یکسان خروجی‌ها تمرکز دارد. این تفاوت، جایگاه پروژه حاضر را در کنار سامانه‌های موجود مشخص می‌کند.

\section{جمع‌بندی}

در این فصل، ابتدا مفاهیم اصلی مورد نیاز برای درک پروژه شامل تحلیل شطرنج، موتورهای شطرنج، مدل‌های زبانی، تولید توضیحات شطرنجی و ارزیابی خروجی مدل‌ها بررسی شد. همچنین مشخص شد که ترکیب توانایی تحلیلی موتورهای شطرنج با قابلیت تولید زبان طبیعی مدل‌های زبانی می‌تواند راهی برای تبدیل تحلیل‌های تخصصی به توضیحات قابل فهم برای انسان باشد. بااین‌حال، مطالعات مرتبط نشان می‌دهند که روانی متن به‌تنهایی تضمین‌کننده صحت توضیح نیست و ارزیابی محتوای شطرنجی اهمیت زیادی دارد.

در ادامه، پژوهش‌های مرتبط از جمله روش‌های تولید توضیح بر اساس داده‌های بازی، ترکیب مدل زبانی با موتور شطرنج، تولید توضیحات هدایت‌شده بر اساس مفاهیم و چارچوب‌های جدید ارزیابی بررسی شدند. این بررسی نشان می‌دهد که اگرچه بخش‌های مختلف مسئله در پژوهش‌های پیشین مورد توجه قرار گرفته‌اند، همچنان امکان بررسی و مقایسه مدل‌های زبانی مختلف در یک مجموعه آزمون یکسان و کنترل‌شده وجود دارد. بر همین اساس، در فصل بعد روش پیشنهادی پروژه، ساختار نرم‌افزار، نحوه تهیه مجموعه آزمون، فرایند تحلیل موقعیت‌ها، تولید توضیحات و روش ارزیابی مدل‌های مختلف تشریح می‌شود.
